\begin{abstract}
We measure whether the same parallel Rust/Rayon code scales the same way
natively and as WebAssembly (WASM) with threads in Chrome and Firefox. On an
i9-14900KF with 8 dedicated P-cores we ran 5310 independent runs of an image
filter (no synchronization) and of an image histogram with shared atomic or
per-task local counters, under four contention levels, and 21\,840 more runs
in nineteen exploratory campaigns on the same build. With enough work per
parallel region, WASM threads scale like native code: on the filter, Chrome's
speedup is 0.996 times native's (99.5\% CI [0.987, 1.004]) and Firefox's
1.055 times [1.034, 1.072], at 1.17 to 1.36 times the native time, a factor
that does not grow with the threads. Hardware counters and SIMD variants
indicate that the filter is bound by memory at eight threads, so Firefox's slower loop scales better; it
does not run threads better. Synchronization depends on the engine: relative
to the lock-free control, shared atomics cost Chrome 1.4 to 2.3 times what
they cost natively under concentrated contention, and 0.42 times when
contention is spread over 256 counters; in Firefox they never cost more.
The penalty follows one instruction sequence: for an atomic add, V8's
optimizing compiler emits a compare-and-swap loop where LLVM and
SpiderMonkey emit one locked add. With V8's baseline compiler the penalty
disappears, and a native program with that loop gives Chrome's one-thread
times within 3\% and both directions of the effect; why the loop wins under
spread contention is not shown. This report is the
long version of the paper: it gives every campaign, probe and pilot of the
study, each decision with the number behind it, every result with its
uncertainty, the problems of browser benchmarking that we met, and what
stays open.
\end{abstract}
